% Figure: energy accounting for a domestic refrigerator over a day. Left bar:
% where the daily heat load into the cabinet comes from (wall conduction,
% door openings, internal loads). Right bar: where the electrical input goes
% (compressor isentropic work, compressor+motor losses, throttling loss,
% condenser/evaporator pinch loss). Simple stacked horizontal bars, no axis env.
\begin{figure}[htbp]
\centering
\begin{tikzpicture}[x=1cm, y=1cm]
% ==== common geometry ====
% bars are 8 cm wide, representing 100 percent.
\def\barW{8}
\def\barH{0.7}
% ===== top bar: heat load into the cabinet (100% = total daily load) =====
% wall conduction 55, door openings 25, internal/defrost/gasket 20
\fill[engblue]      (0,3)      rectangle ({0.55*\barW},{3+\barH});
\fill[engamber]     ({0.55*\barW},3) rectangle ({0.80*\barW},{3+\barH});
\fill[enggray]      ({0.80*\barW},3) rectangle (\barW,{3+\barH});
\node[font=\scriptsize, anchor=east] at (-0.15,{3+\barH/2}) {heat load in};
\node[font=\scriptsize, white] at ({0.275*\barW},{3+\barH/2}) {wall 55\%};
\node[font=\scriptsize, white] at ({0.675*\barW},{3+\barH/2}) {door 25\%};
\node[font=\scriptsize, white] at ({0.90*\barW},{3+\barH/2}) {20\%};
% ===== bottom bar: electrical input destination =====
% reversible (Carnot) minimum 22, compressor+motor loss 33,
% throttling loss 17, heat-exchanger pinch loss 28
\fill[englime!70!black] (0,1.6)      rectangle ({0.22*\barW},{1.6+\barH});
\fill[engred]           ({0.22*\barW},1.6) rectangle ({0.55*\barW},{1.6+\barH});
\fill[engorange]        ({0.55*\barW},1.6) rectangle ({0.72*\barW},{1.6+\barH});
\fill[engblue]          ({0.72*\barW},1.6) rectangle (\barW,{1.6+\barH});
\node[font=\scriptsize, anchor=east] at (-0.15,{1.6+\barH/2}) {electricity in};
\node[font=\scriptsize, white] at ({0.11*\barW},{1.6+\barH/2}) {min 22\%};
\node[font=\scriptsize, white] at ({0.385*\barW},{1.6+\barH/2}) {compr 33\%};
\node[font=\scriptsize, white] at ({0.635*\barW},{1.6+\barH/2}) {throt 17\%};
\node[font=\scriptsize, white] at ({0.86*\barW},{1.6+\barH/2}) {pinch 28\%};
\end{tikzpicture}
\caption[Domestic-refrigerator energy accounting]{Energy accounting for a small
domestic refrigerator over a day of steady operation. The upper bar splits the
heat that leaks into the cold cabinet: conduction through the insulated walls
dominates, with door openings and internal sources (lights, the defrost heater,
gasket infiltration) making up the rest. The lower bar splits the electrical
input the compressor draws to pump that heat back out: only about a fifth is the
reversible minimum the Carnot bound sets for the cabinet-to-room lift, while the
compressor and its motor, the throttling step, and the finite-temperature pinch
in the condenser and evaporator each claim a comparable share. The teardown
shows where a design retrofit buys the most: thicker wall insulation cuts the
load the whole machine must answer, while a larger condenser narrows the pinch
loss, and the two together outrank any plausible gain from the compressor alone.}
\label{fig:vol04ch03:fridge-teardown}
\end{figure}
